\documentclass[11pt]{article}
\usepackage[OT1]{fontenc}
\renewcommand{\rmdefault}{cmr}
\renewcommand{\sfdefault}{cmss}
\renewcommand{\ttdefault}{cmtt}
\usepackage{amsmath,amssymb,amsthm}
\usepackage[hidelinks]{hyperref}
\setlength{\textwidth}{6.1in}
\setlength{\oddsidemargin}{0.2in}
\setlength{\evensidemargin}{0.2in}
\setlength{\textheight}{9in}
\setlength{\topmargin}{-0.35in}
\setlength{\headheight}{13.6pt}
\setlength{\emergencystretch}{1.5em}
\newtheorem{theorem}{Theorem}[section]
\newtheorem{lemma}[theorem]{Lemma}
\newtheorem{proposition}[theorem]{Proposition}
\newtheorem{corollary}[theorem]{Corollary}
\theoremstyle{definition}
\newtheorem{definition}[theorem]{Definition}
\newtheorem{example}[theorem]{Example}
\newcommand{\Comp}{\operatorname{Comp}}
\newcommand{\Forb}{\operatorname{Forb}_{\rm ind}}
\newcommand{\E}{\mathsf E}
\newcommand{\C}{\mathcal C}
\newcommand{\B}{\mathcal B}
\newcommand{\OO}{\mathcal O}
\newcommand{\NN}{\operatorname{N}}
\newcommand{\dSigma}{d\text{-}\Sigma}

\newcommand{\lh}{\operatorname{lh}}
\newcommand{\NF}{\operatorname{NF}}
\newcommand{\wt}{\operatorname{W}}
\newcommand{\tri}{\operatorname{T}}
\title{A calculator for pure first-order resemblance:\\
Cantor normal forms, finite patterns, and Lean proofs}
\author{}
\date{23 October 2026\\Working manuscript}
\begin{document}
\maketitle
\begin{abstract}
Carlson's pure first-order resemblance relation on the ordinals has an
explicit arithmetic description below $\varepsilon_0$. We present a
standalone calculator for that description, together with finite-pattern
operations and a Lean development. The central operation computes the
largest $\beta$ such that $\alpha\leq_1\beta$, using hereditary Cantor
normal forms with compressed natural coefficients. We explain why the
recursion terminates, why repeated summands must be combined by ordinal
rather than commutative arithmetic, and how its result determines finite
traces of the global relation. We also give a sharp quadratic bound on
expanded output-tree size. The Lean development proves arithmetic normality and denotation,
finite-pattern laws, the sharp size bound, and the identification of the
computed endpoint with the independently constructed recursive
$\Sigma_1$-elementarity relation. The semantic proof uses complete
finite diagrams, exact weighted-forest capacity, constructive copies,
and finite obstructions, without assuming Carlson's characterization. The Python
implementation is separately tested against expanded reference
calculations. Finite pattern enumeration supplies syntax and map tests,
not an order-two ordinal assignment or a computation of a transfinite
order type.
\end{abstract}

\section{What the calculator computes}

Two different orders occur in this subject. The usual ordinal order
compares the positions of ordinals. A resemblance relation asks whether
an earlier initial segment has enough of the finite configurations of a
later one. Carlson defines the pure structure
\[
 \mathcal R_1=(\mathrm{Ord};\leq,\leq_1)
\]
by recursion: $\alpha\leq_1\beta$ means that $\alpha\leq\beta$ and
the restriction to $\alpha$ is a $\Sigma_1$-elementary substructure of
the restriction to $\beta$ in this same relational language
\cite{Carlson1999}. The reference to $\leq_1$ inside the language is
part of the recursive definition. It is not ordinary elementarity in
the language of linear order alone.

For a fixed $\alpha$, its $\leq_1$-successors form an interval. Write
$\lh(\alpha)$ for its largest element when that element exists.
Below $\varepsilon_0$ there is a finite arithmetic recipe for this
endpoint. Here $\varepsilon_0$ is the least fixed point of
$\xi\mapsto\omega^\xi$, and a ``finite recipe'' means a terminating
calculation on a finite ordinal notation, not that the ordinal denoted
is finite. The result decides the relation by
\begin{equation}\label{eq:interval}
 \alpha\leq_1\beta
 \quad\Longleftrightarrow\quad
 \alpha\leq\beta\leq\lh(\alpha).
\end{equation}
The semantic characterization is Carlson's Theorem~2.5 of
\cite{Carlson1999}. Section~\ref{sec:semantic} explains the constructive
argument formalized here for the calculator's range. Wilken presents the endpoint recursion
in a convenient Cantor-normal-form formulation in Proposition~1.1 of
\cite{Wilken2021}.

The distinction matters even in small examples. The calculator gives
$\lh(\omega)=\omega+1$, so $\omega\leq_1\omega+1$, whereas
$\omega\not\leq_1\omega+2$. It does not merely test whether the two
notations are in increasing order.

This paper explains an implementation of those published results.
The recursion and its semantic identification are credited to their
sources. The output-size proof, representation choices, finite map
interface, and explicit division between tested and formalized claims
are given here so that the package can be used without another project.

\section{Hereditary Cantor normal forms}

Every nonzero ordinal below $\varepsilon_0$ has a unique expression
\[
 a=\omega^{e_1}c_1+\cdots+\omega^{e_k}c_k,
 \qquad e_1>\cdots>e_k,\qquad 0<c_i<\omega,
\]
where each exponent is itself represented in the same way. Zero is the
empty expression. The Python representation is a tuple of pairs
$(e_i,c_i)$, recursively using tuples for the exponents. Natural
coefficients are ordinary integers and are stored without expanding
$c_i$ copies of the monomial.

A valid normal form has positive coefficients, recursively normal
exponents, and strictly decreasing exponents at the outer level.
These conditions are checked on programmatic inputs. For example, two
separate consecutive $\omega^2$ terms are not a canonical tuple; their
canonical form combines their coefficients. The readable input layer
constructs normal forms using the same arithmetic operations.

\subsection{Addition and finite right multiplication}

Ordinal addition is essential here. If the leading exponent of $b$ is
$e$, adding $b$ to $a$ removes the trailing monomials of $a$ with
exponent less than $e$, combines a remaining equal-exponent monomial
if one exists, and then appends the remainder of $b$. Thus
\[
 1+\omega=\omega,\qquad \omega+1>\omega.
\]
The recursive normal-form comparison is lexicographic: compare leading
exponents, then their coefficients, then the tails. On valid hereditary
Cantor normal forms this is the usual ordinal order.

For a positive natural $c$, the operation $a\cdot c$ means repeated
addition on the right. If $a=\omega^e d+t>0$ is in normal form, then
\[
 a\cdot c=\omega^e(dc)+t.
\]
Only the final copy retains the lower tail. In particular,
\begin{equation}\label{eq:noncommutative}
 (\omega+1)\cdot2=\omega\cdot2+1.
\end{equation}
Treating the coefficients as a reason to commute or distribute all
summands would change the endpoint calculation.

\section{The compressed endpoint recursion}

Let $\operatorname{last}(a)$ be the final exponent of a nonzero normal
form, and put $\operatorname{last}(0)=0$. Define an auxiliary function
$P$ by
\begin{align}
 P(0)&=0,\label{eq:pzero}\\
 P(\omega^e c+t)&=(\omega^e+P(e))\cdot c+P(t).
 \label{eq:precursion}
\end{align}
Define the executable endpoint function by
\begin{equation}\label{eq:h}
 H(0)=0,\qquad H(a)=a+P(\operatorname{last}(a))\quad(a>0).
\end{equation}
All calls in \eqref{eq:precursion} enter an exponent subtree or a proper
tail of the input. The recursion therefore terminates structurally.
Neither its termination argument nor its definition consults a
resemblance oracle.

\subsection{Why this is Wilken's formula}

An additive-principal ordinal is a power of $\omega$. Expand an exponent
$e$ as its weakly decreasing list of additive-principal summands
$\rho_1,\ldots,\rho_m$, including repetitions from coefficients.
For $\rho_i=\omega^{d_i}$, its computed endpoint is
$\rho_i+P(d_i)$. Consequently \eqref{eq:precursion} gives
\begin{equation}\label{eq:expanded}
 P(e)=H(\rho_1)+\cdots+H(\rho_m).
\end{equation}
The order of this sum is retained. Compression simply replaces a block
of identical successive summands by finite right multiplication.
Substituting \eqref{eq:expanded} into \eqref{eq:h} gives exactly the
recursion in \cite[Proposition~1.1]{Wilken2021}. The semantic proof in
Section~\ref{sec:semantic} identifies $H(a)$ with $\lh(a)$, giving
\eqref{eq:interval} and reconstructing the relevant part of Carlson's theorem.

There are two correctness steps here. The first proves that a compressed
algorithm denotes the displayed expanded ordinal sum. The second
identifies that arithmetic sum with the recursively defined semantic
relation. A successful program execution establishes neither theorem
by itself. Both steps have general Lean proofs in this package;
Section~\ref{sec:lean} describes their exact statements and dependencies.

\subsection{Worked examples}

\begin{center}
\begin{tabular}{c|c}
Input $a$ & Endpoint $H(a)$\\\hline
$0$ & $0$\\
$7$ & $7$\\
$\omega$ & $\omega+1$\\
$\omega^2$ & $\omega^2+2$\\
$\omega^\omega$ & $\omega^\omega+\omega+1$\\
$\omega^{\omega+1}$ & $\omega^{\omega+1}+\omega+2$\\
$\omega^{\omega\cdot2}$ & $\omega^{\omega\cdot2}+\omega\cdot2+1$
\end{tabular}
\end{center}

For the last example, the last exponent is $\omega\cdot2$. Its two
principal summands are both $\omega$, whose endpoints are $\omega+1$.
Their sum is \eqref{eq:noncommutative}, which explains the final $+1$,
rather than $+2$. This example tests the coefficient compression at a
place where ordinary commutative arithmetic would give the wrong answer.

For a successor ordinal $a$, the final exponent is zero, so $H(a)=a$.
For nonzero $e$, the recursion shows inductively that $P(e)$ has a
positive finite final summand. Hence the endpoint of a nonzero limit
below $\varepsilon_0$ is a successor. In particular,
\[
 H(H(a))=H(a).
\]
This is a useful consistency check on the arithmetic routine, although
it is not a replacement for its denotation proof.

\section{Finite patterns and maps}

The second part of the package operates on finite relational syntax.
Fix the underlying linear order $0<\cdots<n-1$. A pure pattern of order
$k\geq1$ is represented by strict relations $R_1,\ldots,R_k$ on this
set. Reflexive pairs are implicit. Starting with $R_0$ equal to the
strict linear order, require at each level:
\begin{enumerate}
 \item $R_i\subseteq R_{i-1}$;
 \item $R_i$ is transitive;
 \item if $xR_{i-1}yR_{i-1}z$ and $xR_i z$, then $xR_i y$.
\end{enumerate}
These are the finite nesting and respect conditions used in the pure
pattern setting; see the introduction of \cite{Wilken2021}. The
implementation checks them on the finite relations. Its order-two
objects remain finite syntax: constructing one does not assign its
least ordinal realization.

\subsection{Coverings preserve; induced embeddings also reflect}

A covering from $A$ to $B$ is an increasing injection $f$ that preserves
every higher relation:
\[
 xR_i^A y\ \Longrightarrow\ f(x)R_i^B f(y).
\]
An induced embedding additionally reflects the relations, replacing the
implication by an equivalence. On the fixed finite linear order, all
increasing injections can be enumerated by their increasing image lists.
Thus a search through the $\binom{|B|}{|A|}$ possible image sets is a
complete finite map search, with no candidate when $|A|>|B|$.

For example, let $A$ be a two-point pattern with no strict $R_1$ edge,
and let $B$ have the edge $0R_1 1$. The identity on the underlying
two-point order is a covering $A\to B$: there are no positive $R_1$
facts in $A$ to violate. It is not an induced embedding, because it
does not reflect the new edge of $B$. Conversely, there is no covering
$B\to A$.

Both notions have identities and compose. Increasing injections compose,
and relation preservation passes through the intermediate pattern.
For induced embeddings the two reflection implications compose as well.
These elementary laws are useful checks on an implementation because
confusing preservation with reflection changes the research problem.

\subsection{Markers and finite traces}

A pattern may carry a distinguished point. It need not be the maximum.
A map is required to preserve that point only when marker preservation
is explicitly requested; that request requires both markers to exist.
The number of higher relations determines the pattern's order, regardless
of where its marker lies.

For strictly increasing normal-form coordinates $a_0<\cdots<a_{n-1}$,
the calculator can form the finite trace whose strict first relation is
\[
 iR_1j\quad\Longleftrightarrow\quad
 i<j\ \hbox{ and }\ a_j\leq H(a_i).
\]
This is the restriction of the global pure relation, using the
semantic characterization proved in Section~\ref{sec:semantic}. It is not obtained by recomputing elementarity inside
the finite list itself. For instance the trace on
$\{\omega,\omega+1\}$ retains a genuine edge of $\mathcal R_1$.
Recomputing a recursive elementary-substructure relation on a finite
linear order gives a different object, since a new element beyond a
finite last parameter is witnessed by an existential formula.

\section{What the finite enumeration counts}

The implementation enumerates the strict relations satisfying the stated
laws on the fixed underlying order. Up to five points it gives:
\begin{center}
\begin{tabular}{r|rrrrrr}
Number of points & $0$ & $1$ & $2$ & $3$ & $4$ & $5$\\\hline
Order one & $1$ & $1$ & $2$ & $5$ & $14$ & $42$\\
Order two & $1$ & $1$ & $3$ & $14$ & $84$ & $596$
\end{tabular}
\end{center}
The empty pattern is included. Fixing the linear order eliminates
nontrivial automorphisms, so these are finite syntactic isomorphism-type
counts. They are not counts of covering-equivalence classes, minimal
ordinal realizations, or ordinal ranks. In particular, their numerical
growth cannot verify a conjecture about a transfinite order type.

The search is exhaustive only for its supplied finite size. Its
combinatorial cost grows quickly; the command-line interface therefore
keeps large enumerations behind an explicit size limit. No empirical
failure to find a pattern or a large covering is used as an unbounded
negative theorem.

\section{A sharp expanded-tree size bound}\label{sec:size}

Compact storage size and the mathematical norm of a notation are
different quantities. Define the expanded hereditary monomial-tree
weight by
\begin{equation}\label{eq:weight}
 \wt(0)=0,\qquad
 \wt(\omega^e c+t)=c(1+\wt(e))+\wt(t).
\end{equation}
Repeated exponent subtrees are charged repeatedly, and coefficients are
unary for this measure. Put $\tri(n)=n(n+1)/2$.

\begin{theorem}\label{thm:size}
For every normal form $a$ below $\varepsilon_0$,
\[
 \wt(P(a))\leq\tri(\wt(a)),\qquad
 \wt(H(a))\leq\tri(\wt(a)).
\]
The second bound is attained for every natural input weight.
\end{theorem}

\begin{proof}
Ordinal addition deletes a final segment and possibly combines two
equal-exponent coefficient blocks. It introduces no new exponent
subtree. Therefore
\[
 \wt(u+v)\leq\wt(u)+\wt(v),\qquad
 \wt(u\cdot c)\leq c\wt(u).
\]
Also $\tri(x)+\tri(y)\leq\tri(x+y)$ and
$c\tri(x)\leq\tri(cx)$ for natural $c,x,y$.

For the first inequality, use structural induction. Zero is immediate.
Write $a=\omega^e c+t$, $r=\wt(e)$, and $s=\wt(t)$. The
recursion and induction hypotheses give
\begin{align*}
 \wt(P(a))
 &\leq c(1+r+\wt(P(e)))+\wt(P(t))\\
 &\leq c(1+r+\tri(r))+\tri(s)\\
 &=c\tri(r+1)+\tri(s)\\
 &\leq\tri(c(r+1)+s)=\tri(\wt(a)).
\end{align*}
For nonzero $a$ put $n=\wt(a)$ and $e=\operatorname{last}(a)$.
At least one monomial of $a$ contains that exponent, so
$1+\wt(e)\leq n$. Hence
\[
 \wt(H(a))\leq n+\wt(P(e))
 \leq n+\tri(n-1)=\tri(n).
\]
At zero the same assertion is equality.

For sharpness define $t_0=0$ and $t_{n+1}=\omega^{t_n}$.
Then $\wt(t_n)=n$, and direct induction in the recursion gives
\[
 H(t_n)=t_n+t_{n-1}+\cdots+t_1\quad(n\geq1).
\]
The displayed summands are in decreasing Cantor normal form. Their
expanded weights sum to $n+(n-1)+\cdots+1=\tri(n)$. Use zero
for $n=0$.
\end{proof}

This theorem concerns output size under \eqref{eq:weight}. It is not a
quadratic running-time claim for arbitrary Python integers, a bound in
binary coefficient length, or a theorem comparing ordinal notations
with minimum pattern size. Those require separate cost models and
mathematical arguments.

\section{A finite-forest proof of the semantic endpoint}\label{sec:semantic}

We now explain the mathematical bridge from the arithmetic computation to
Carlson's recursively defined relation. This bridge is proved in Lean in
this package. The proof follows the constructive route suggested in the
concluding section of \cite{Carlson1999}, using weighted finite forests.
It does not import a correctness assumption for the endpoint operation.

Write $A(x,y)$ for the arithmetic candidate
$x\leq y\leq H(x)$ when $x$ has a normal notation. In a bounded initial
segment every actual ordinal has such a notation: the formalization proves
this by well-founded induction using the zero, successor, and cofinal
limit cases of Mathlib's fundamental sequences. Thus notation syntax does
not silently replace quantification over actual ordinal elements.

Several arithmetic laws prepare the argument. On normal inputs, $P$
preserves ordinal addition and is strictly increasing. Both $P$ and $H$
preserve every strict additive-principal bound. In addition,
\[
 y\leq P(e)\Longrightarrow H(y)\leq P(e),\qquad
 x\leq y\leq H(x)\Longrightarrow H(y)\leq H(x).
\]
These are proved from the recursion, independently of elementarity. The
second law makes the intervals of $A$ into a forest: intervals issued
from comparable points either nest or separate. Finite restrictions are
therefore finite ordered forests of exactly the kind used below.

\subsection{Weights and exact finite capacity}

Regard a finite ordered forest as an empty forest or a first root with
children forest $C$, followed by a remaining forest $R$. Define its
weight by
\[
 w(\emptyset)=0,\qquad
 w(\operatorname{cons}(C,R))=\omega^{w(C)}+w(R).
\]
The sum is taken in the displayed order and then normalized. For example,
a leaf followed by a one-child tree has weight $1+\omega=\omega$.
This absorption is compatible with finite realization: the extra leaf
fits before the first infinite root.

A greedy placement after a cut $\sigma$ puts the first root at
$\sigma+\omega^{w(C)}$, places $C$ after that root, and begins $R$
after the root's computed endpoint. It is an induced realization of the
whole forest, including its negative relation facts. Its final extent is
\begin{equation}\label{eq:capacity}
 \sigma+P(w(F)).
\end{equation}
For a nonempty forest this is its greatest label. For an empty forest,
``extent'' means the cut itself.

\begin{lemma}\label{lem:capacity}
Suppose an induced arithmetic realization of $F$ has every label strictly
after $\sigma$ and at most $b$, where $\sigma\leq b$. Then
$\sigma+P(w(F))\leq b$. The greedy placement attains equality.
\end{lemma}

\begin{proof}
Induct on the forest. If its first root is $r$, the induction hypothesis
on the children gives
$r+P(w(C))\leq H(r)=r+P(\operatorname{last}(r))$.
Strict monotonicity of $P$ yields
$w(C)\leq\operatorname{last}(r)$. Since
$\omega^{\operatorname{last}(r)}$ divides $r$, ordinal rounding gives
$\sigma+\omega^{w(C)}\leq r$ whenever $\sigma<r$.
The first greedy root and its endpoint therefore precede or equal the
corresponding bounds in the given realization. Every point of the
remaining forest is strictly beyond $H(r)$: its negative relation to
$r$ is part of being an induced realization. Apply the induction
hypothesis to that remaining forest with cut $H(r)$. Additivity of $P$
combines the two bounds into \eqref{eq:capacity}. Empty children and an
empty remainder contribute zero. The same recursion computes the greedy
extent and proves that all its positive and negative atoms agree with
the forest.
\end{proof}

\subsection{Copies inside the endpoint interval}

Suppose $\alpha<\beta\leq H(\alpha)$, and let a finite configuration
below $\beta$ be given. Its old points lie below $\alpha$; call its new
points $Y$. Write the final principal block as
$\alpha=\xi+\omega^e$, with $\xi<\alpha$.
By finite-forest extraction, $Y$ has an induced forest description $F$.
Every point of $Y$ is after $\xi$ and strictly below
\[
 H(\alpha)=\xi+P(\omega^e).
\]
Lemma~\ref{lem:capacity} and strict monotonicity of $P$ imply
$w(F)<\omega^e$. Consequently $P(w(F))<\omega^e$ as well.

Choose a cut $c<\alpha$ above $\xi$, all old points, and the endpoints
$H(x)$ of the old points for which $A(x,\alpha)$ is false. There are
only finitely many such requirements, and each is below $\alpha$.
The greedy copy after $c$ ends at $c+P(w(F))<\alpha$, since both
its remaining displacement and the remainder of $c$ lie in the final
additive-principal block. It has exactly the induced diagram of $Y$.

The choice of $c$ also pays every mixed old/new atom. If $A(x,\alpha)$
holds, the copied new points remain in $x$'s endpoint interval. Otherwise
they lie beyond $H(x)$. Equality and order are preserved by keeping the
old points fixed and putting the new copy strictly above them. This gives
one induced copy of the whole finite configuration below $\alpha$.
The equal-endpoint case uses the identity map.

\subsection{A finite obstruction beyond the endpoint}

Every normal notation $e$ has a canonical finite forest of weight $e$:
expand its coefficients into repeated roots and recurse into its
exponents. Put that forest under one new root. Placed greedily after
$\xi$, its root is $\alpha=\xi+\omega^e$, and all its labels are
at most $H(\alpha)$. Include $\xi$ as an old fixed parameter.

An induced copy of this finite diagram below $\alpha$ would send its
root to some $r$ with $\xi<r<\alpha$. The children still have weight
$e$, so the root-capacity step in Lemma~\ref{lem:capacity} gives
$e\leq\operatorname{last}(r)$. Rounding then implies
$\alpha=\xi+\omega^e\leq r$, a contradiction. At $\alpha=0$,
a one-point diagram supplies the obstruction. Thus no point strictly
beyond $H(\alpha)$ can satisfy the arithmetic finite-copy criterion.

\subsection{Identification with recursive elementarity}

The package proves Carlson's finite-set criterion
\cite[Proposition~1.2]{Carlson1999} for the explicit single-relation
language: existential Boolean-matrix truth reflects exactly when every
finite configuration has an induced copy fixing its old parameters.
Complete finite diagrams record negative atoms as well as positive ones.
This proof includes empty domains and empty tuples.

The copy construction and obstruction show that $A$ satisfies the
$\Sigma_1$ initial-segment equation below every normal-notation bound.
A separate upper-ordinal induction proves uniqueness of the recursively
constructed relation satisfying that equation there. Hence $A$ agrees
with the actual pure $\mathcal R_1$ relation on normal inputs.
Downwardness and the obstruction at $H(\alpha)+1$ extend the result to
every actual right ordinal:
\[
 \alpha\leq_1\beta
 \quad\Longleftrightarrow\quad
 \alpha\leq\beta\leq H(\alpha).
\]
In particular, $H(\alpha)$ is proved to be the greatest endpoint, rather
than merely a correct answer among a finite selection of notations.
The formal statement is restricted on the left to the calculator's
hereditary normal forms below $\varepsilon_0$; the periodic extension
beyond that range is outside this calculator.

\section{Lean development and scope}\label{sec:lean}

The Lean package pins its toolchain and Mathlib revision \cite{Mathlib}.
Its arithmetic uses Mathlib's \texttt{ONote} syntax and normal subtype
\texttt{NONote}. Positive coefficients are typed as positive natural
numbers; normal inputs carry proof obligations. The recursive semantics
uses the precise signature $(=,\leq,\leq_1)$ and Boolean matrices with
arbitrary finite existential witness blocks, matching Carlson's stated
meaning of $\Sigma_1$.

The principal semantic declaration is
\texttt{PureR1.endpoint\_semantic}. Given any normal input $a$ and
any actual ordinal $\beta$, it proves
\[
 \operatorname{R1}(a.\operatorname{repr},\beta)
 \quad\Longleftrightarrow\quad
 a.\operatorname{repr}\leq\beta\leq H(a).\operatorname{repr}.
\]
The relation on the left is constructed by well-founded upper-ordinal
recursion from formula truth. It is not defined as the endpoint interval.
The declarations \texttt{endpoint\_isGreatest} and
\texttt{intervalCheck\_iff\_sigma1} give the greatest-endpoint and
executable-checker consequences. Their types contain no supplied
characterization, reflection oracle, or semantic correctness premise.

The semantic proof includes both directions of the finite-copy criterion,
coverage of actual ordinal cuts by normal notations, full induced forest
extraction and realization, the finite copying construction and endpoint
obstruction, and recursive uniqueness. The norm development proves
Theorem~\ref{thm:size}, its triangular formula, and matching tower
witnesses for every natural input size. The finite-pattern development
proves the stated checker and map laws, distinguishing coverings from
induced embeddings and making marker preservation explicit.

Selected formula, recursive-prefix, and finite-diagram constructions were
adapted from an earlier Alexander thesis Lean development. Their
single-relation versions are included and checked locally. The larger
project is not a dependency; its simultaneous order-two relations are
not substituted for the standalone relation used here.

The proofs use Lean and Mathlib's standard foundations, with principal
axiom reports limited to \texttt{propext}, \texttt{Classical.choice}
and \texttt{Quot.sound}. No admitted proofs, added axioms or external
Carlson characterization are used. The coverage guide lists the actual
modules and declarations, and the verification record identifies their
checked source hashes.

Python and Lean remain different implementations. Python's parser,
interpreter and execution are tested separately; there is no formal
cross-language refinement proof. Lean's finite existence search enumerates
finite functions, whereas Python enumerates increasing image lists.
The finite enumeration table reports Python executions. Neither the
checker laws nor the norm theorem asserts a running-time bound, a minimum
order-two realization, or an extension to additive or simultaneous
higher-order resemblance.

\section{Using and reproducing the package}

The project contains the Python package, standard-library tests, a
command-line interface, the Lean source and Lake configuration, this
manuscript, and a local PDF build script. Its readable ordinal input is
a small specified grammar; it is parsed as data rather than evaluated
as Python code. Machine-readable output includes a versioned format.
The precise command syntax and examples are in \texttt{USAGE.md}.

The basic Python check is
\begin{verbatim}
python3 -m unittest discover -s tests -v
\end{verbatim}
For a fresh Lean installation, use the toolchain and dependency revisions
specified in the \texttt{lean/} directory. After acquiring the pinned
dependencies and their compiled cache, run the project build. The
optional \texttt{scripts/check\_lean.py} compiles named modules serially
with explicit time and memory limits; it does not update dependencies.
Its locally borrowed cache is a development convenience, not a
repository dependency.

The complete local check is \texttt{scripts/verify.py --lean --paper}.
It compiles all local Lean modules in dependency order, checks the stated
axiom dependencies, runs the Python suite and explicit shared fixtures,
and rebuilds the PDF. The records identify source hashes and actual
commands. A resource refusal remains an incomplete computation rather
than a negative mathematical answer.

\section{Research role and directions for reuse}

The calculator was developed as a first-order base case for work on
finite patterns and ordinal notations. A motivating question in
Carlson's earlier cohabitating-trees manuscript concerns the computability
and order types of patterns ordered by their least ordinal coverings
\cite[Section~5]{CarlsonDraft}. That author draft conjectures that the
order-two type $\kappa_2$ is the proof-theoretic ordinal of
$\mathrm{KP}\ell_0$. The finite order-two syntax interface here is
preparation for such questions, not their ordinal evaluator. Related
marked pure-order-two notation results are already supplied by
\cite{Wilken2018}; the formulation of the source question must be
checked before treating it as open.

Three immediate uses are concrete. First, endpoint and finite-map
calculations provide regression cases for more elaborate pattern
algorithms. Second, the expanded weight and Theorem~\ref{thm:size}
support controlled-size calculations, where changing the coefficient
encoding can change the statement. Third, finite traces supply explicit
examples separating inherited resemblance from a relation recomputed
inside a finite structure. A catalogue of finite certificate tasks can
combine those operations without claiming to measure an agent's
intrinsic proof-theoretic ordinal.

Extensions need to respect the language. Pure $\mathcal R_1$, additive
patterns, and the first relation inside simultaneous higher-order
structures are different semantic objects. Wilken's introduction
\cite{Wilken2021} explains where the simple pure repetition ceases to
describe the first relation of $\mathcal R_2$. Reusing the Cantor-normal-
form arithmetic is often possible; reusing the resemblance decision
without a transfer theorem is not. A complete order-two ordinal
assignment, a higher-pattern comparison procedure, and a Bachmann--Howard
termination-order implementation each require additional mathematics.

\begin{thebibliography}{9}
\bibitem{Carlson1999}
T. J. Carlson, \emph{Ordinal arithmetic and $\Sigma_1$-elementarity},
Archive for Mathematical Logic \textbf{38} (1999), 449--460.

\bibitem{Wilken2021}
G. Wilken, \emph{Pure $\Sigma_2$-elementarity beyond the core},
Annals of Pure and Applied Logic \textbf{172} (2021), article 103001.
\url{https://arxiv.org/abs/1710.01870v5}.

\bibitem{Wilken2018}
G. Wilken, \emph{Pure patterns of order 2},
Annals of Pure and Applied Logic \textbf{169} (2018), 54--82.
\url{https://arxiv.org/abs/1608.08421}.

\bibitem{CarlsonDraft}
T. J. Carlson, \emph{Generalizing Kruskal's theorem to pairs of
cohabitating trees}, author manuscript dated 6 October 2008, Section~5.
Published version: Archive for Mathematical Logic \textbf{55} (2016),
37--48, \url{https://doi.org/10.1007/s00153-015-0457-4}.
The concluding conjecture passage cited here is in the earlier author
manuscript and is not reproduced in the final journal body.

\bibitem{Mathlib}
The mathlib Community, \emph{Mathlib4},
\href{https://github.com/leanprover-community/mathlib4/tree/360da6fa66c1273b76b6b2d8c5666fd5ac2e3b56}{source revision \texttt{360da6fa}},
especially \texttt{Mathlib.SetTheory.Ordinal.Notation}.
\end{thebibliography}
\end{document}
